\section{Controle de Atitude}
O controle de atitude de uma espaçonave consiste em projetar leis de realimentação capazes de estabilizar e orientar o corpo em torno de uma referência desejada.
Nesta dissertação será empregado o controle proporcional-derivativo (PD), pela sua implementação direta e eficiência local de estabilização.

%==================================================
%----- Erro de atitude e PD
%==================================================
\subsection{Erro de atitude e lei de controle PD}
Sejam os quaternions unitários para a atitude atual e para a atitude desejada, respectivamente:
\begin{equation}
q = [q_0,q_x,q_y,q_z,]^T,
\end{equation}

\begin{equation}
q_{ref} = [q_{0(ref)},q_{x(ref)},q_{y(ref)},q_{z(ref)}]^T,
\end{equation}

Define-se o quaternion de erro de atitude como:
\begin{equation}
q_e = q_{\mathrm{ref}}^{-1} \otimes q =
[\,q_{e(0)},q_{e(x)},q_{e(y)},q_{e(z)}\,]^T,
\label{eq:definicao_quaternion_erro}
\end{equation}
onde $\otimes$ é o produto de quaternions, e $[\,q_{e0},\vec\epsilon_e\,]^T$ suas partes escalar e vetorial, análogo a \eqref{eq:def_quaternion}.
Para evitar que ocorra esforço de controle ineficiente adota-se o vetor de erro:
\begin{equation}
\vec e_q = 2\,sgn(q_{e0})\,\vec\epsilon_e,
\label{eq:erro_quaternion_vetorial}
\end{equation}

em que $sgn(\cdot)$ é a função sinal, podendo se tornar um dos casos a seguir:

\begin{equation}
\operatorname{sgn}(x) =
\begin{cases}
+1, & x > 0, \\
0,  & x = 0, \\
-1, & x < 0,
\end{cases}
\label{eq:casos_funcao_sinal}
\end{equation}

Como $q_e$ e seu negativo ($-q_e$) representam a mesma orientação física, existe um recobrimento extra na representação.
O termo $sgn(q_{e})$ faz com que a parte escalar do quaternion de erro seja escolhida de forma consistente ($q_{e}\ge 0$), evitando o fenômeno conhecido como \emph{unwinding}, no qual o controlador tende a forçar uma rotação desnecessária de $360^{\circ}$ para atingir a referência.
Portanto, o vetor de erro $\vec e_q$ em \eqref{eq:erro_quaternion_vetorial} fornece uma medida proporcional ao menor ângulo de rotação necessário para alinhar a base em relação à atitude de referência.



%==================================================
%----- Arquitetura geral
%==================================================
\section{Arquitetura de controle base--manipulador}
A arquitetura de controle adotada para a operação do manipulador em microgravidade contém três malhas principais: (i) o controle da atitude da base, (ii) o controle do manipulador com mínima reação, e (iii) o controle por torque nas juntas. Essas malhas operam de forma coordenada.

\subsection{Controle de atitude da base}
A primeira malha de controle é responsável por controlar a orientação da base do sistema.
Neste trabalho foi adotada a convenção de Hamilton\fn{Detalhado anteriormente em \eqref{sec:Quaternions}}, assim, o erro de quaternion é dado por:
\begin{equation}
    q_e \;=\; q_{ref}^{-1} \otimes q_b \;=\;
    \begin{bmatrix}
        q_{e} \\[2pt] \vec {\epsilon}_e
    \end{bmatrix},
    \label{eq:erro_quaternion1}
\end{equation}

onde:
\begin{itemize}
    \item $q_{e}$ é a parte escalar do quaternion de erro e $\vec {\epsilon}_e$ é o vetor dos eixos associados;
    \item $q_b$ é o quaternion da atitude da base;
    \item $q_{\mathrm{ref}}$ é o quaternion de referência e $q_{ref}^{-1}$ o seu inverso;
    \item $\vec e_q$ define o vetor de erro de atitude usado na realimentação.
\end{itemize}


Caso a referência apresente velocidade angular não nula, define-se também a velocidade de erro na mesma referência, como:
\begin{equation}
    \vec{\omega}_e \;=\; \vec{\omega}_b \;-\; \vec{\omega}_{\mathrm{ref}},
    \label{eq:erro_velocidade}
\end{equation}

em que $\vec{\omega}_b$ é a velocidade angular da base e $\vec{\omega}_{\mathrm{ref}}$ a velocidade angular de referência.
Com base nessas definições, a lei de controle proporcional-derivativa (PD) aplicada à base é

\begin{equation}
    \vec{\tau}_b=
    -\,K_p\,\vec e_q
    -K_d\,\vec\omega_e,
    \label{eq:lei_controle_PD}
\end{equation}

onde $K_p$ e $K_d$ são os ganhos proporcionais e derivativos, respectivamente. 
Neste trabalho o termo integral não é utilizado, de modo a evitar o acúmulo de erro ao longo do tempo (wind--up) 
e por não ser usualmente empregado no ambiente espacial.




\section{Velocidade angular}
A velocidade angular de referência $\vec \omega_{ref}$ está no \emph{frame} analisado, podendo ser qualquer elo ou outra parte do manipulador robótico.
Para compará-lo com a velocidade angular $\vec{\omega_b}$ da base, mapeia-se $\vec{\omega}_{ref}$ para o frame através da rotação associada ao erro:
\begin{equation}
\vec{\omega}_e \;=\; \vec{\omega} \;-\; \mathbf{R}(q_e)\,\vec{\omega}_{ref}.
\label{eq:omega_e}
\end{equation}
onde $\vec \omega_e$ é o erro entre as velocidades angulares, $\vec \omega$ é a velocidade angular medida e
$\vec \omega_{ref}$ a desejada.
Com essa definição, a cinemática do erro é dada por:
\begin{equation}
\dot{q}_e \;=\; \tfrac{1}{2}\,\Omega(q_e)\,\vec{\omega}_e,
\label{eq:qedot}
\end{equation}
onde $\Omega(\cdot)$ é o mapeamento apresentado em \eqref{eq:definicao_qdot}.
Para manter a consistência numérica, normaliza-se periodicamente $q_e \leftarrow \frac{q_e}{\|q_e\|}$.

A lei de controle PD para a atitude é dada por:
\begin{equation}
\vec{M} \;=\; -\,\mathbf{K}_p\,\vec{e}_q \;-\; \mathbf{K}_d\,\vec{\omega}_e,
\label{eq:controle_PD_fundamental}
\end{equation}
com $\mathbf{K}_p={}^T\mathbf{K}_p \succ 0$ e $\mathbf{K}_d={}^T\mathbf{K}_d \succ 0$.

%==================================================
%----- Caso particular
%==================================================
Para o caso particular onde $\vec{\omega}_{ref}=\vec{0}$, e
$\vec{\omega}_e=\vec{\omega}$ poderá ser consultado no Apêndice \ref{ApendiceC}.

%==================================================
%----- PD veloc angular
%==================================================
\subsection{PD aplicado às velocidades angulares}
O controlador PD atua sobre a velocidade angular da espaçonave, buscando minimizar o erro entre a atitude desejada e a atual.
Seja $\vec{\omega}_0$ a velocidade angular medida e $\vec{\omega}_1$ a velocidade angular desejada (referência).
De forma similar a \eqref{eq:controle_PD_fundamental}, o controle de torque de controle pode ser expresso matematicamente como:
\begin{equation}
\vec{M}_{\text{PD}} = -K_p \, \vec{e}_q - K_d \, (\vec{\omega}_0 - \vec{\omega}_1),
\label{eq:torque_PD}
\end{equation}

Onde $K_p$ e $K_d$ são os ganhos escalares de controle, e $\vec{e}_q$ representa o erro de atitude, que pode ser expresso em termos de quaternions ou ângulos de Euler linearizados.  
Nesta dissertação, o erro de atitude será tratado em termos de quaternions, de modo a evitar singularidades associadas aos ângulos de Euler.